\documentclass[aspectratio=169, 10pt]{beamer}
\usepackage{sunass_estilo}

\title{Módulo I: Teoría del Consumidor}
\subtitle{Guía Teórica de Repaso Rápido · PEU SUNASS}
\author[PEU SUNASS · Economía]{Preparación Examen de Admisión PEU SUNASS}
\institute{Basado en el Temario Oficial de Economía · Microeconomía}
\date[2026]{2026}

\begin{document}

% ==============================================================================
% SLIDE 1: PORTADA
% ==============================================================================
\begin{frame}[plain]
  \titlepage
\end{frame}

% ==============================================================================
% SLIDE 2: ESTRUCTURA DEL MÓDULO (TARJETAS + RECUADRO)
% ==============================================================================
\begin{frame}{Mapa Conceptual del Módulo I}
  \vspace{0.3cm}
  \begin{columns}[T]
    \begin{column}{0.32\textwidth}
      \tarjeta{azulPrincipal}{1. Restricción Presupuestaria}{
        Lo que el consumidor \textbf{puede} comprar dados precios \((p_1, p_2)\) e ingreso \(m\).\\
        \(p_1 x_1 + p_2 x_2 \le m\)
      }
    \end{column}
    \begin{column}{0.32\textwidth}
      \tarjeta{naranjaAcento}{2. Preferencias y Utilidad}{
        Lo que el consumidor \textbf{desea} comprar. Axiomas, curvas de indiferencia y \(TMS\).
      }
    \end{column}
    \begin{column}{0.32\textwidth}
      \tarjeta{verdeAcento}{3. Elección y Demanda}{
        Maximización de \(U(x_1, x_2)\) sujeta al presupuesto: demanda marshalliana y Slutsky.
      }
    \end{column}
  \end{columns}

  \vspace{0.4cm}
  \recuadroclave{
    \textbf{Problema Central:} \(\max_{x_1, x_2} U(x_1, x_2)\) \quad sujeto a \quad \(p_1 x_1 + p_2 x_2 = m\)
  }

  \vspace{0.2cm}
  \mensajefinal{El equilibrio del consumidor iguala su valoración subjetiva (\(TMS\)) con el costo de oportunidad del mercado (\(p_1/p_2\)).}
\end{frame}

% ==============================================================================
% SLIDE 3: RESTRICCIÓN PRESUPUESTARIA (GRÁFICO IZQ + CIFRAS DER)
% ==============================================================================
\begin{frame}{1. Restricción Presupuestaria y Quiebres}
  \begin{columns}[T]
    \begin{column}{0.54\textwidth}
      {\small\textbf{Conjunto Presupuestario y Desplazamientos}}
      \vspace{0.1cm}
      \begin{tikzpicture}[scale=0.82]
        \begin{axis}[
          axis lines=left,
          xlabel={$x_1$}, ylabel={$x_2$},
          xmin=0, xmax=12, ymin=0, ymax=10,
          xtick={6, 10}, xticklabels={$\frac{m}{p_1'}$, $\frac{m}{p_1}$},
          ytick={7}, yticklabels={$\frac{m}{p_2}$},
          width=8.2cm, height=5.4cm,
          clip=false
        ]
          % Recta inicial
          \addplot[aceroAcento, very thick, domain=0:10] {7 - 0.7*x};
          \node[aceroAcento, font=\scriptsize, anchor=west] at (axis cs:6.2, 3.2) {Inicial: $-\frac{p_1}{p_2}$};
          % Recta con aumento de p1
          \addplot[rojoAcento, very thick, dashed, domain=0:6] {7 - 1.1667*x};
          \node[rojoAcento, font=\scriptsize, anchor=east] at (axis cs:3.2, 2.2) {$p_1' > p_1$};
          % Area sombreada
          \addplot[fill=azulClaroFondo, opacity=0.45, draw=none] coordinates {(0,0) (0,7) (6,0)} \closedcycle;
        \end{axis}
      \end{tikzpicture}
      {\scriptsize\color{grisTexto}\textit{Nota: Si todos los precios y el ingreso se multiplican por \(\lambda > 0\), la recta no cambia.}}
    \end{column}

    \begin{column}{0.43\textwidth}
      \indicador{Ecuación de la Recta}{aceroAcento}{\(x_2 = \frac{m}{p_2} - \frac{p_1}{p_2}x_1\)}
      \indicador{Pendiente (Costo de Oportunidad)}{rojoAcento}{\(-\frac{p_1}{p_2}\)}
      \indicador{Impuesto ad-valorem vs. específico}{naranjaAcento}{\(p_1(1+\tau)\) \quad vs. \quad \(p_1 + t\)}
      {\scriptsize
        \textbf{Regla rápida para quiebres (SUNASS):}\\
        • \textbf{Bloques tarifarios crecientes:} si \(p_1\) sube a partir de \(\bar{x}_1\), la recta se quiebra hacia adentro (más inclinada).\\
        • \textbf{Subsidio en especie / bono:} desplazamiento paralelo horizontal con tramo plano.
      }
    \end{column}
  \end{columns}

  \mensajefinal{Un aumento proporcional de todos los precios equivalente a una caída del ingreso real mantiene intacto el poder de compra relativo.}
\end{frame}

% ==============================================================================
% SLIDE 4: PREFERENCIAS, AXIOMAS Y TMS
% ==============================================================================
\begin{frame}{2. Preferencias, Utilidad y Tasa Marginal de Sustitución}
  \vspace{0.15cm}
  \begin{columns}[T]
    \begin{column}{0.32\textwidth}
      \tarjeta{azulPrincipal}{Axiomas Racionales}{
        \textbf{1. Completitud:} toda cesta es comparable.\\
        \textbf{2. Transitividad:} \(A \succeq B, B \succeq C \Rightarrow A \succeq C\).\\
        \textbf{3. Monotonicidad:} ``más es mejor'' (pendiente negativa).\\
        \textbf{4. Convexidad:} se prefieren promedios.
      }
    \end{column}
    \begin{column}{0.32\textwidth}
      \tarjeta{naranjaAcento}{Función de Utilidad}{
        Es un concepto \textbf{ordinal}, no cardinal.\\[4pt]
        Cualquier transformación monótona creciente \(V = f(U)\) con \(f'>0\) representa las \textbf{mismas preferencias} y deja inalterada la \(TMS\).
      }
    \end{column}
    \begin{column}{0.32\textwidth}
      \tarjeta{verdeAcento}{Tasa Marginal de Sustitución}{
        Cantidad de \(x_2\) a la que se renuncia por \(+1\) unidad de \(x_1\) manteniendo \(U\) constante:\\[4pt]
        \(TMS_{1,2} = -\frac{dx_2}{dx_1}\Big|_{dU=0} = \frac{UMg_1}{UMg_2}\)
      }
    \end{column}
  \end{columns}

  \vspace{0.1cm}
  \recuadroclave{
    \small\textbf{Propiedad clave:} Las curvas de indiferencia \textbf{nunca se cruzan} (violaría transitividad + monotonicidad).
  }
  \mensajefinal{La \(TMS\) decreciente equivale geométricamente a curvas de indiferencia estrictamente convexas al origen.}
\end{frame}

% ==============================================================================
% SLIDE 5: TIPOS DE PREFERENCIAS ESPECIALES (GRÁFICO + FÓRMULAS)
% ==============================================================================
\begin{frame}{2b. Funciones de Utilidad Clásicas en Exámenes}
  \begin{columns}[T]
    \begin{column}{0.52\textwidth}
      {\small\textbf{Mapa Comparativo de Curvas de Indiferencia}}
      \vspace{0.05cm}
      \begin{tikzpicture}[scale=0.72]
        % Panel 1: Cobb-Douglas
        \draw[->, thick] (0,0) -- (3.2,0) node[right, font=\scriptsize] {$x_1$};
        \draw[->, thick] (0,0) -- (0,2.8) node[above, font=\scriptsize] {$x_2$};
        \draw[azulPrincipal, very thick, domain=0.45:2.8, samples=40] plot (\x, {1.2/\x});
        \node[font=\scriptsize, azulPrincipal] at (1.6, 2.3) {\textbf{Cobb-Douglas}};
        \node[font=\tiny, grisTexto] at (1.6, 1.9) {$U = x_1^\alpha x_2^\beta$};

        % Panel 2: Sustitutos Perfectos
        \begin{scope}[xshift=4.0cm]
          \draw[->, thick] (0,0) -- (3.2,0) node[right, font=\scriptsize] {$x_1$};
          \draw[->, thick] (0,0) -- (0,2.8) node[above, font=\scriptsize] {$x_2$};
          \draw[naranjaAcento, very thick] (0,2.2) -- (2.2,0);
          \draw[naranjaAcento, thick, dashed] (0,1.4) -- (1.4,0);
          \node[font=\scriptsize, naranjaAcento] at (1.7, 2.3) {\textbf{Sust. Perfectos}};
          \node[font=\tiny, grisTexto] at (1.7, 1.9) {$U = ax_1 + bx_2$};
        \end{scope}

        % Panel 3: Complementos Perfectos
        \begin{scope}[yshift=-3.4cm]
          \draw[->, thick] (0,0) -- (3.2,0) node[right, font=\scriptsize] {$x_1$};
          \draw[->, thick] (0,0) -- (0,2.6) node[above, font=\scriptsize] {$x_2$};
          \draw[verdeAcento, very thick] (0.8, 2.3) -- (0.8, 0.8) -- (2.8, 0.8);
          \draw[grisTexto, dotted] (0,0) -- (2.5, 2.5);
          \node[font=\scriptsize, verdeAcento] at (2.1, 2.1) {\textbf{Leontief}};
          \node[font=\tiny, grisTexto] at (2.1, 1.7) {$U = \min\{ax_1, bx_2\}$};
        \end{scope}

        % Panel 4: Cuasilineal
        \begin{scope}[xshift=4.0cm, yshift=-3.4cm]
          \draw[->, thick] (0,0) -- (3.2,0) node[right, font=\scriptsize] {$x_1$};
          \draw[->, thick] (0,0) -- (0,2.6) node[above, font=\scriptsize] {$x_2$};
          \draw[rojoAcento, very thick, domain=0.2:2.6, samples=30] plot (\x, {2.0 - 0.8*sqrt(\x)});
          \draw[rojoAcento, thick, dashed, domain=0.2:2.6, samples=30] plot (\x, {2.5 - 0.8*sqrt(\x)});
          \node[font=\scriptsize, rojoAcento] at (1.9, 2.3) {\textbf{Cuasilineal}};
          \node[font=\tiny, grisTexto] at (1.9, 1.9) {$U = v(x_1) + x_2$};
        \end{scope}
      \end{tikzpicture}
    \end{column}

    \begin{column}{0.45\textwidth}
      \indicador{Cobb-Douglas: Gasto Fijo}{azulPrincipal}{\(x_1^* = \frac{\alpha}{\alpha+\beta}\frac{m}{p_1}\)}
      \indicador{Sustitutos Perfectos}{naranjaAcento}{Solución de Esquina si \(\frac{a}{b} \neq \frac{p_1}{p_2}\)}
      \indicador{Complementos (Leontief)}{verdeAcento}{\(x_1^* = \frac{m}{p_1 + \frac{a}{b}p_2}\) \ (Vértice: \(ax_1 = bx_2\))}
      {\scriptsize
        \textbf{¡Ojo en la Cuasilineal!}\\
        La demanda de \(x_1\) \textbf{no depende del ingreso} \(m\) en soluciones interiores (\(\text{Efecto Ingreso} = 0\)). Todo aumento de ingreso va a \(x_2\).
      }
    \end{column}
  \end{columns}

  \mensajefinal{Identificar la forma funcional al instante ahorra el 70\% del tiempo de cálculo en el examen.}
\end{frame}

% ==============================================================================
% SLIDE 6: ELECCIÓN ÓPTIMA DEL CONSUMIDOR (TANGENCIA VS ESQUINA)
% ==============================================================================
\begin{frame}{3. Elección Óptima: Lagrangiano y Soluciones de Esquina}
  \begin{columns}[T]
    \begin{column}{0.53\textwidth}
      {\small\textbf{Equilibrio Interior (Tangencia)}}
      \vspace{0.1cm}
      \begin{tikzpicture}[scale=0.82]
        \begin{axis}[
          axis lines=left,
          xlabel={$x_1$}, ylabel={$x_2$},
          xmin=0, xmax=10, ymin=0, ymax=9,
          xtick={4.5}, xticklabels={$x_1^*$},
          ytick={4.5}, yticklabels={$x_2^*$},
          width=8.0cm, height=5.3cm
        ]
          % Recta presupuestaria
          \addplot[aceroAcento, very thick, domain=0:9] {9 - x};
          % Curvas de indiferencia
          \addplot[grisTexto, dashed, domain=1.2:8.5, samples=40] {12/x};
          \addplot[azulPrincipal, very thick, domain=2.3:8.8, samples=50] {20.25/x};
          \addplot[rojoAcento, dotted, thick, domain=3.5:8.8, samples=40] {30/x};
          % Punto óptimo
          \addplot[only marks, mark=*, mark size=2.8pt, color=verdeAcento] coordinates {(4.5, 4.5)};
          \draw[dotted, thick, verdeAcento] (axis cs:4.5,0) -- (axis cs:4.5,4.5) -- (axis cs:0,4.5);
          \node[verdeAcento, font=\scriptsize, anchor=south west] at (axis cs:4.6, 4.6) {\textbf{Óptimo \(E^*\)}};
        \end{axis}
      \end{tikzpicture}
    \end{column}

    \begin{column}{0.44\textwidth}
      \indicador{Condición de Primer Orden (Interior)}{azulPrincipal}{\(TMS = \frac{UMg_1}{UMg_2} = \frac{p_1}{p_2}\)}
      \indicador{Principio Equimarginal}{verdeAcento}{\(\frac{UMg_1}{p_1} = \frac{UMg_2}{p_2} = \lambda\)}
      {\small
        \textbf{¿Cuándo hay Solución de Esquina?}\\
        Si en \((0, m/p_2)\) se cumple \(TMS < \frac{p_1}{p_2}\), el consumidor gasta todo en \(x_2\): \(x_1^* = 0\).\\
        \vspace{0.2cm}
        \textbf{Interpretación de \(\lambda\):}\\
        Utilidad marginal del ingreso \(\left(\frac{\partial V}{\partial m}\right)\).
      }
    \end{column}
  \end{columns}

  \mensajefinal{El Lagrangiano falla si la solución es de esquina o si las preferencias son no convexas (cóncavas).}
\end{frame}

% ==============================================================================
% SLIDE 7: DESCOMPOSICIÓN DE SLUTSKY (GRÁFICO CLÁSICO)
% ==============================================================================
\begin{frame}{4. Demanda Individual: Ecuación de Slutsky}
  \begin{columns}[T]
    \begin{column}{0.53\textwidth}
      {\small\textbf{Caída de \(p_1\): Efecto Sustitución y Efecto Ingreso}}
      \vspace{0.1cm}
      \begin{tikzpicture}[scale=0.82]
        \begin{axis}[
          axis lines=left,
          xlabel={$x_1$}, ylabel={$x_2$},
          xmin=0, xmax=12, ymin=0, ymax=9,
          xtick={2.5, 4.8, 7.2}, xticklabels={$x_1^A$, $x_1^S$, $x_1^B$},
          ytick=\empty,
          width=8.2cm, height=5.4cm
        ]
          % Recta inicial A
          \addplot[aceroAcento, very thick, domain=0:5.5] {7.5 - 1.36*x};
          % Recta nueva B (p1 baja)
          \addplot[verdeAcento, very thick, domain=0:11] {7.5 - 0.68*x};
          % Recta compensada Slutsky (pasa por A con nueva pendiente)
          \addplot[naranjaAcento, thick, dashed, domain=0:8.5] {5.8 - 0.68*x};

          % Puntos A, S, B
          \addplot[only marks, mark=*, mark size=2.5pt, aceroAcento] coordinates {(2.5, 4.1)};
          \node[aceroAcento, font=\scriptsize, anchor=south west] at (axis cs:2.4,4.2) {$A$};
          \addplot[only marks, mark=*, mark size=2.5pt, naranjaAcento] coordinates {(4.8, 2.53)};
          \node[naranjaAcento, font=\scriptsize, anchor=north east] at (axis cs:4.8,2.5) {$S$};
          \addplot[only marks, mark=*, mark size=2.5pt, verdeAcento] coordinates {(7.2, 2.6)};
          \node[verdeAcento, font=\scriptsize, anchor=south west] at (axis cs:7.2,2.6) {$B$};

          % Flechas de efectos
          \draw[->, thick, naranjaAcento] (axis cs:2.5, 0.7) -- (axis cs:4.8, 0.7) node[midway, above, font=\tiny] {ES};
          \draw[->, thick, verdeAcento] (axis cs:4.8, 0.7) -- (axis cs:7.2, 0.7) node[midway, above, font=\tiny] {EI};
        \end{axis}
      \end{tikzpicture}
    \end{column}

    \begin{column}{0.44\textwidth}
      \indicador{Ecuación de Slutsky}{azulPrincipal}{\(\frac{\partial x_1}{\partial p_1} = \frac{\partial x_1^h}{\partial p_1} - x_1 \frac{\partial x_1}{\partial m}\)}
      {\scriptsize
        \textbf{Signos Invariantes (Ley de la Demanda):}\\
        • \textbf{Efecto Sustitución (ES):} Siempre negativo (\(\le 0\)). Si baja \(p_1\), siempre sube \(x_1\) compensado.\\
        • \textbf{Bien Normal:} EI refuerza al ES (\(\frac{\partial x_1}{\partial m} > 0\)) \(\Rightarrow\) demanda decreciente.\\
        • \textbf{Bien Inferior:} EI va en contra del ES.\\
        • \textbf{Bien Giffen:} Bien \textit{muy inferior} donde \(|EI| > |ES|\) \(\Rightarrow\) curva de demanda con pendiente positiva.\\[4pt]
        \textbf{Identidad de Roy:} \(x_i(p,m) = -\frac{\partial V/\partial p_i}{\partial V/\partial m}\)
      }
    \end{column}
  \end{columns}

  \mensajefinal{Todo bien Giffen es necesariamente inferior, pero no todo bien inferior es un bien Giffen.}
\end{frame}

% ==============================================================================
% SLIDE 8: DEMANDA DE MERCADO, ELASTICIDADES Y APLICACIÓN SUNASS
% ==============================================================================
\begin{frame}{5. Demanda de Mercado, Elasticidades y Saneamiento}
  \vspace{0.1cm}
  \begin{columns}[T]
    \begin{column}{0.32\textwidth}
      \tarjeta{azulPrincipal}{Agregación de Mercado}{
        \(X_1(p_1) = \sum_{i=1}^n x_1^i(p_1, p_2, m_i)\)\\[4pt]
        Suma \textbf{horizontal} de cantidades a cada precio (bienes privados).\\
        \textit{Cuidado con tramos donde algunos consumidores salen del mercado (quiebres).}
      }
    \end{column}
    \begin{column}{0.32\textwidth}
      \tarjeta{naranjaAcento}{Elasticidad y Gasto Total}{
        \(\varepsilon_p = \frac{\partial Q}{\partial P}\frac{P}{Q}\)\\[4pt]
        • \(|\varepsilon_p| > 1\): Elástica (\(\downarrow P \Rightarrow \uparrow \text{Ingreso}\)).\\
        • \(|\varepsilon_p| < 1\): Inelástica (\(\uparrow P \Rightarrow \uparrow \text{Ingreso}\)).\\
        • \(\varepsilon_m > 0\): Normal (\(\varepsilon_m > 1\): lujo).
      }
    \end{column}
    \begin{column}{0.32\textwidth}
      \tarjeta{verdeAcento}{Aplicación SUNASS}{
        El agua potable residencial es un bien \textbf{necesario e inelástico} (\(|\varepsilon_p| \approx 0.1 - 0.3\)).\\[4pt]
        Justifica regulación tarifaria para evitar abusos monopólicos y diseñar \textbf{subsidios cruzados}.
      }
    \end{column}
  \end{columns}

  \vspace{0.05cm}
  \recuadroclave{
    \small\textbf{Dualidad Clave:} Marshalliana \(x(p,m) \leftrightarrow V(p,m)\) \quad | \quad Hicksiana \(h(p,u) \leftrightarrow e(p,u)\) (Shephard: \(h_i = \frac{\partial e}{\partial p_i}\))
  }
  \mensajefinal{En bienes inelásticos como el agua potable, un alza tarifaria eleva la recaudación pero reduce el excedente del consumidor.}
\end{frame}

% ==============================================================================
% SLIDE 9: CURVA DE ENGEL, SEI Y PREFERENCIAS HOMOTÉTICAS
% ==============================================================================
\begin{frame}{6. Curvas de Engel, SEI y Preferencias Homotéticas}
  \begin{columns}[T]
    \begin{column}{0.52\textwidth}
      {\small\textbf{Curvas de Engel (\(m\) vs. \(x_1\))}}
      \vspace{0.05cm}
      \begin{tikzpicture}[scale=0.78]
        \begin{axis}[
          axis lines=left,
          xlabel={$x_1$}, ylabel={Ingreso ($m$)},
          xmin=0, xmax=10, ymin=0, ymax=9,
          xtick=\empty, ytick=\empty,
          width=8.0cm, height=5.2cm
        ]
          % Homotética / Proporcional
          \addplot[azulPrincipal, very thick, domain=0:8.5] {x};
          \node[azulPrincipal, font=\scriptsize, anchor=south east] at (axis cs:8.4, 8.2) {Homotética ($\varepsilon_m=1$)};
          % Necesario (cóncava en x vs m -> m crece rápido, x lento; en plano (x1, m) es convexa hacia eje m)
          \addplot[verdeAcento, very thick, domain=0:6.5] {0.2*x^2};
          \node[verdeAcento, font=\scriptsize, anchor=west] at (axis cs:6.5, 8.1) {Necesario ($\varepsilon_m<1$)};
          % Lujo
          \addplot[naranjaAcento, very thick, domain=0:8.8, samples=40] {2.5*sqrt(x)};
          \node[naranjaAcento, font=\scriptsize, anchor=north west] at (axis cs:5.8, 5.7) {Lujo ($\varepsilon_m>1$)};
          % Inferior
          \addplot[rojoAcento, very thick, dashed, domain=0:6] {8 - 0.8*(x-3)^2};
          \node[rojoAcento, font=\scriptsize, anchor=west] at (axis cs:4.8, 3.2) {Inferior ($\varepsilon_m<0$)};
        \end{axis}
      \end{tikzpicture}
    \end{column}

    \begin{column}{0.45\textwidth}
      \indicador{Preferencias Homotéticas}{azulPrincipal}{\(SEI = \text{Rayo lineal desde el origen}\)}
      {\scriptsize
        • La \(TMS\) solo depende del cociente \(x_2/x_1\), \textbf{no de la escala}.\\
        • Ejemplos: Cobb-Douglas, Leontief, CES, Sustitutos perfectos.\\
        • Todas las elasticidades ingreso son \(\varepsilon_m = 1\).\\[6pt]
      }
      \indicador{Clasificación por \(\varepsilon_m\)}{verdeAcento}{\(\varepsilon_m > 1 \text{ Lujo} \ \mid \ 0<\varepsilon_m<1 \text{ Necesario}\)}
      {\scriptsize
        • \textbf{Senda de Expansión del Ingreso (SEI):} une las cestas óptimas al variar \(m\) manteniendo \(p_1, p_2\) fijos.\\
        • \textbf{Curva Precio-Consumo (CPC):} une óptimos al variar \(p_1\).
      }
    \end{column}
  \end{columns}

  \mensajefinal{No todos los bienes pueden ser inferiores ni todos de lujo al mismo tiempo: su promedio ponderado es siempre \(1\).}
\end{frame}

% ==============================================================================
% SLIDE 10: PREFERENCIA REVELADA (WARP, SARP E ÍNDICES)
% ==============================================================================
\begin{frame}{7. Preferencia Revelada (WARP / SARP) e Índices}
  \begin{columns}[T]
    \begin{column}{0.52\textwidth}
      {\small\textbf{Detección Gráfica del Axioma Débil (WARP)}}
      \vspace{0.05cm}
      \begin{tikzpicture}[scale=0.78]
        \begin{axis}[
          axis lines=left,
          xlabel={$x_1$}, ylabel={$x_2$},
          xmin=0, xmax=10, ymin=0, ymax=9,
          xtick=\empty, ytick=\empty,
          width=8.0cm, height=5.2cm
        ]
          % Recta en t=0
          \addplot[aceroAcento, very thick, domain=0:9] {6.5 - 0.72*x};
          % Recta en t=1
          \addplot[rojoAcento, very thick, domain=0:6.5] {8.5 - 1.3*x};
          % Cestas A y B dentro de la intersección -> Violación WARP
          \addplot[only marks, mark=*, mark size=2.8pt, aceroAcento] coordinates {(3.5, 3.98)};
          \node[aceroAcento, font=\scriptsize, anchor=south west] at (axis cs:3.6, 4.0) {\textbf{\(A\) (en \(t=0\))}};
          \addplot[only marks, mark=*, mark size=2.8pt, rojoAcento] coordinates {(2.2, 5.64)};
          \node[rojoAcento, font=\scriptsize, anchor=south west] at (axis cs:2.3, 5.6) {\textbf{\(B\) (en \(t=1\))}};
          % Punto C consistente
          \addplot[only marks, mark=square*, mark size=2.5pt, verdeAcento] coordinates {(7.0, 1.46)};
          \node[verdeAcento, font=\scriptsize, anchor=south west] at (axis cs:7.0, 1.5) {\(C\)};
        \end{axis}
      \end{tikzpicture}
      {\scriptsize\color{grisTexto}\textit{\(A\) y \(B\) son asequibles en ambos periodos: elegir \(A\) y luego \(B\) viola WARP.}}
    \end{column}

    \begin{column}{0.45\textwidth}
      \indicador{Axioma Débil (WARP)}{rojoAcento}{\(p^0 x^0 \ge p^0 x^1 \implies p^1 x^1 < p^1 x^0\)}
      {\scriptsize
        • Si \(A\) fue elegida cuando \(B\) era asequible, entonces \(B\) \textbf{nunca} puede ser elegida cuando \(A\) es asequible.\\
        • \textbf{SARP (Fuerte):} extiende WARP a cadenas transitivas (condición necesaria y suficiente para maximizar utilidad).\\[5pt]
      }
      \indicador{Índices de Bienestar}{verdeAcento}{\(L_q = \frac{p^1 x^1}{p^1 x^0} \quad \text{vs.} \quad P_q = \frac{p^0 x^1}{p^0 x^0}\)}
      {\scriptsize
        • Si en \(t=1\) puede comprar \(x^0\) (\(p^1 x^1 > p^1 x^0\)), está \textbf{mejor}.\\
        • Si en \(t=0\) podía comprar \(x^1\) (\(p^0 x^0 > p^0 x^1\)), está \textbf{peor}.
      }
    \end{column}
  \end{columns}

  \mensajefinal{WARP permite contrastar la racionalidad del consumidor usando solo precios y cantidades observadas.}
\end{frame}

% ==============================================================================
% SLIDE 11: MEDIDAS DE BIENESTAR (VC, VE Y EXCEDENTE DEL CONSUMIDOR)
% ==============================================================================
\begin{frame}{8. Medidas de Bienestar: \(VC\), \(VE\) y \(\Delta EC\)}
  \begin{columns}[T]
    \begin{column}{0.52\textwidth}
      {\small\textbf{Alza de Precio (\(p_1^0 \to p_1^1\)) para un Bien Normal}}
      \vspace{0.05cm}
      \begin{tikzpicture}[scale=0.78]
        \begin{axis}[
          axis lines=left,
          xlabel={$x_1$}, ylabel={$p_1$},
          xmin=0, xmax=10, ymin=0, ymax=8,
          xtick=\empty,
          ytick={2.5, 5.5}, yticklabels={$p_1^0$, $p_1^1$},
          width=8.0cm, height=5.2cm
        ]
          % Demanda Hicksiana U1 (izquierda)
          \addplot[verdeAcento, very thick, dashed, domain=1.5:7.5] {8.5 - 0.9*x};
          \node[verdeAcento, font=\tiny, anchor=south] at (axis cs:6.4, 2.7) {$h(U^1)$};
          % Demanda Marshalliana (centro)
          \addplot[azulPrincipal, very thick, domain=1.5:8.5] {9.2 - 0.85*x};
          \node[azulPrincipal, font=\tiny, anchor=south] at (axis cs:7.6, 2.8) {$x(p,m)$};
          % Demanda Hicksiana U0 (derecha)
          \addplot[rojoAcento, very thick, dotted, domain=2.5:9.5] {10.2 - 0.85*x};
          \node[rojoAcento, font=\tiny, anchor=south] at (axis cs:8.8, 2.8) {$h(U^0)$};

          % Lineas de precios
          \draw[grisTexto, dashed] (axis cs:0, 2.5) -- (axis cs:9.1, 2.5);
          \draw[grisTexto, dashed] (axis cs:0, 5.5) -- (axis cs:5.6, 5.5);
        \end{axis}
      \end{tikzpicture}
    \end{column}

    \begin{column}{0.45\textwidth}
      \indicador{Variación Compensatoria (\(VC\))}{rojoAcento}{\(VC = e(p^1, U^0) - e(p^0, U^0)\)}
      {\scriptsize
        Dinero necesario tras el cambio para volver a la utilidad \textbf{inicial \(U^0\)}.\\[4pt]
      }
      \indicador{Variación Equivalente (\(VE\))}{verdeAcento}{\(VE = e(p^1, U^1) - e(p^0, U^1)\)}
      {\scriptsize
        Dinero que habría que quitar antes del cambio para dejarlo en la utilidad \textbf{final \(U^1\)}.\\[5pt]
      }
      {\scriptsize
        \textbf{Orden de Magnitud (Bien Normal):}\\
        • Si el precio \textbf{sube}: \(|VE| < |\Delta EC| < |VC|\).\\
        • Si el precio \textbf{baja}: \(|VC| < |\Delta EC| < |VE|\).\\
        • \textbf{Cuasilineal (\(EI=0\)):} \(VC = \Delta EC = VE\).
      }
    \end{column}
  \end{columns}

  \mensajefinal{El Excedente del Consumidor (\(\Delta EC\)) solo es una medida exacta del bienestar cuando no hay efecto ingreso.}
\end{frame}

% ==============================================================================
% SLIDE 12: PROPIEDADES DE LA DEMANDA Y SLUTSKY VS HICKS
% ==============================================================================
\begin{frame}{9. Propiedades de la Demanda y Slutsky vs. Hicks}
  \vspace{0.1cm}
  \begin{columns}[T]
    \begin{column}{0.32\textwidth}
      \tarjeta{azulPrincipal}{Propiedades Marshallianas}{
        \textbf{1. Homogeneidad de grado 0:} \(x(\lambda p, \lambda m) = x(p,m)\) (sin ilusión monetaria).\\
        \textbf{2. Agregación de Engel:} \(\sum w_i \varepsilon_{i,m} = 1\).\\
        \textbf{3. Agregación de Cournot:} \(\sum w_i \varepsilon_{i,j} = -w_j\).
      }
    \end{column}
    \begin{column}{0.32\textwidth}
      \tarjeta{naranjaAcento}{Matriz de Slutsky (\(S_{ij}\))}{
        \(S_{ij} = \frac{\partial h_i}{\partial p_j} = \frac{\partial^2 e}{\partial p_i \partial p_j}\)\\[3pt]
        • \textbf{Simétrica:} \(S_{ij} = S_{ji}\).\\
        • \textbf{Semidefinida negativa:} \(S_{ii} \le 0\) (Ley de la demanda compensada).
      }
    \end{column}
    \begin{column}{0.32\textwidth}
      \tarjeta{verdeAcento}{Brutos vs. Netos y Compensación}{
        • \textbf{Sustitutos brutos:} \(\frac{\partial x_i}{\partial p_j} > 0\) (no simétrico).\\
        • \textbf{Sustitutos netos:} \(\frac{\partial h_i}{\partial p_j} > 0\) (simétrico).\\
        • \textbf{Slutsky vs. Hicks:} Slutsky fija poder de compra (\(x^0\)); Hicks fija utilidad (\(U^0\)).
      }
    \end{column}
  \end{columns}

  \vspace{0.15cm}
  \recuadroclave{
    \small\textbf{Fórmula de Slutsky en Elasticidades:} \quad \(\varepsilon_{i,p_i} = \varepsilon_{i,p_i}^h - w_i \varepsilon_{i,m}\) \quad (donde \(w_i = \frac{p_i x_i}{m}\) es la participación en el gasto)
  }
  \mensajefinal{Si un bien tiene una participación \(w_i\) muy pequeña en el presupuesto, su efecto ingreso es despreciable y \(\Delta EC \approx VC \approx VE\).}
\end{frame}

% ==============================================================================
% SLIDE 13: INCERTIDUMBRE, AVERSIÓN AL RIESGO Y UTILIDAD ESPERADA (vNM)
% ==============================================================================
\begin{frame}{10. Incertidumbre, Aversión al Riesgo y Utilidad Esperada}
  \begin{columns}[T]
    \begin{column}{0.52\textwidth}
      {\small\textbf{Individuo Averso al Riesgo (\(U''(W) < 0\))}}
      \vspace{0.05cm}
      \begin{tikzpicture}[scale=0.78]
        \begin{axis}[
          axis lines=left,
          xlabel={Riqueza ($W$)}, ylabel={Utilidad $U(W)$},
          xmin=0, xmax=10, ymin=0, ymax=3.5,
          xtick={2, 4.5, 5, 8}, xticklabels={$W_1$, $EC_r$, $E[W]$, $W_2$},
          ytick=\empty,
          width=8.2cm, height=5.0cm
        ]
          % U(W) = sqrt(W)
          \addplot[azulPrincipal, very thick, domain=0.2:9.5, samples=50] {sqrt(x)};
          % Cuerda entre (2, 1.414) y (8, 2.828)
          \addplot[grisTexto, dashed, thick, domain=2:8] {1.414 + 0.2357*(x-2)};
          % E[W]=5, U(5)=2.236, E[U]=2.121, EC_r=4.5
          \draw[dotted, rojoAcento, thick] (axis cs:5, 0) -- (axis cs:5, 2.236);
          \draw[dotted, verdeAcento, thick] (axis cs:0, 2.121) -- (axis cs:5, 2.121);
          \draw[dotted, verdeAcento, thick] (axis cs:4.5, 0) -- (axis cs:4.5, 2.121);
          \draw[<->, thick, rojoAcento] (axis cs:4.5, 0.45) -- (axis cs:5.0, 0.45) node[midway, above, font=\tiny] {\(\pi\)};
        \end{axis}
      \end{tikzpicture}
      {\scriptsize\color{grisTexto}\textit{Desigualdad de Jensen: \(U(E[W]) > E[U(W)]\) para \(U\) estrictamente cóncava.}}
    \end{column}

    \begin{column}{0.45\textwidth}
      \indicador{Utilidad Esperada (vNM)}{azulPrincipal}{\(E[U] = p\,U(W_1) + (1-p)\,U(W_2)\)}
      {\scriptsize
        • \textbf{Averso al riesgo:} \(U'' < 0\) (prefiere \(E[W]\) seguro a la lotería).\\
        • \textbf{Neutral:} \(U'' = 0\) (lineal). \quad \textbf{Amante:} \(U'' > 0\) (convexa).\\[4pt]
      }
      \indicador{Equivalente Cierto y Prima de Riesgo}{rojoAcento}{\(U(EC_r) = E[U(W)], \quad \pi = E[W] - EC_r > 0\)}
      {\scriptsize
        • \textbf{Arrow-Pratt:} \(ARA(W) = -\frac{U''(W)}{U'(W)}\) y \(RRA(W) = -W\frac{U''(W)}{U'(W)}\).\\
        • Con seguro actuarialmente justo, el averso al riesgo compra \textbf{cobertura total}.
      }
    \end{column}
  \end{columns}

  \mensajefinal{La función de utilidad vNM es cardinal (única salvo transformaciones afines positivas \(a + bU\) con \(b>0\)), a diferencia de la utilidad ordinal.}
\end{frame}

% ==============================================================================
% SLIDE 14: DOTACIONES INICIALES Y ELECCIÓN INTERTEMPORAL
% ==============================================================================
\begin{frame}{11. Dotaciones Iniciales y Elección Intertemporal}
  \vspace{0.05cm}
  \begin{columns}[T]
    \begin{column}{0.48\textwidth}
      \tarjeta{azulPrincipal}{Slutsky con Dotación Inicial \((\omega_1, \omega_2)\)}{
        Presupuesto: \(p_1 x_1 + p_2 x_2 = p_1 \omega_1 + p_2 \omega_2 = m\).\\[2pt]
        \(\frac{\partial x_1}{\partial p_1} = \underbrace{\frac{\partial x_1^h}{\partial p_1}}_{ES\,(-)} + \underbrace{(\omega_1 - x_1)\frac{\partial x_1}{\partial m}}_{EI \text{ ordinario} + \text{Efecto Dotación}}\)\\[3pt]
        • Si eres \textbf{vendedor neto} (\(\omega_1 > x_1\)) y sube \(p_1\), te vuelves más rico y podrías demandar más \(x_1\) aun siendo bien normal.
      }
    \end{column}
    \begin{column}{0.48\textwidth}
      \tarjeta{verdeAcento}{Elección Intertemporal \((c_1, c_2)\)}{
        Restricción: \(c_1 + \frac{c_2}{1+r} = y_1 + \frac{y_2}{1+r}\) (pendiente \(-(1+r)\)).\\[2pt]
        • Óptimo: \(TMS_{1,2} = \frac{UMg(c_1)}{\beta\, UMg(c_2)} = 1 + r\).\\
        • \textbf{Si sube la tasa de interés \(r\):}\\
        \quad - \textbf{Prestatario} (\(c_1 > y_1\)): empeora y reduce \(c_1\) (ES y EI van en contra de \(c_1\)).\\
        \quad - \textbf{Prestamista} (\(c_1 < y_1\)): mejora su bienestar; efecto ambiguo sobre \(c_1\) (y sobre el ahorro).
      }
    \end{column}
  \end{columns}

  \vspace{0.08cm}
  \recuadroclave{
    \small\textbf{Preferencia Revelada Intertemporal:} Un agente que sigue siendo prestamista tras una subida de \(r\) nunca puede pasar a ser prestatario (violación de WARP); pero un prestatario sí puede volverse prestamista si sube \(r\).
  }
  \mensajefinal{Cuando la riqueza depende de los precios (dotaciones iniciales), un alza del precio del bien vendido puede incrementar el bienestar del agente.}
\end{frame}

% ==============================================================================
% SLIDE 15: CONCLUSIONES / CHEAT-SHEET FINAL DEL MÓDULO I
% ==============================================================================
\begin{frame}{Conclusiones y Formulario Rápido — Módulo I}
  \vspace{0.05cm}
  \recuadroclave{
    \small\textbf{Regla de Oro:} El consumidor optimiza cuando la última unidad monetaria gastada en cada bien rinde exactamente la misma utilidad marginal.
  }

  \vspace{0.15cm}
  \begin{columns}[T]
    \begin{column}{0.32\textwidth}
      \tarjeta{verdeAcento}{Fórmulas Imprescindibles}{
        • \(TMS = \frac{UMg_1}{UMg_2} = \frac{p_1}{p_2}\)\\
        • Cobb-Douglas: \(x_1^* = \frac{\alpha}{\alpha+\beta}\frac{m}{p_1}\)\\
        • Slutsky: \(\Delta x = \Delta x^s + \Delta x^n\)
      }
    \end{column}
    \begin{column}{0.32\textwidth}
      \tarjeta{azulPrincipal}{Trampas de Examen}{
        • Confundir transformación monótona con cambio de preferencias.\\
        • Aplicar tangencia en Leontief o sustitutos perfectos.\\
        • Olvidar el signo negativo en la Identidad de Roy.
      }
    \end{column}
    \begin{column}{0.32\textwidth}
      \tarjeta{rojoAcento}{Siguientes Pasos}{
        • Practicar Bloque B del Módulo 01 en el Simulador (Q31–Q65).\\
        • Revisar curvas de Engel y quiebres tarifarios por bloques.
      }
    \end{column}
  \end{columns}

  \vspace{0.4cm}
  \begin{center}
    {\color{azulPrincipal}\Large\textbf{¡Éxito!}} \quad {\color{grisTexto}Módulo I · Teoría del Consumidor · PEU SUNASS 2026}
  \end{center}
\end{frame}

\end{document}
